% paper/sections/04_grid.tex
\section{界面適合非一様格子の構築}
\label{sec:grid}

\subsection{グリッド密度関数}

界面近傍の解像度を上げるため，Level Set 値に連動したグリッド密度関数 $\omega(\phi)$ を導入する．
\begin{equation}
  \omega(\phi) = 1 + (\alpha - 1)\,\delta_\varepsilon^*(\phi)
\end{equation}
$\alpha > 1$ は界面近傍の格子密度倍率，$\delta_\varepsilon^*(\phi)$ は界面近傍に集中した正規化スムーズデルタ関数である．

\subsubsection{グリッド密度デルタ関数 $\delta_\varepsilon^*(\phi)$ の定義}

$\delta_\varepsilon^*(\phi)$ は界面（$\phi=0$）近傍でのみ大きな値を持ち，
全域積分が正規化条件 $\int_{-\infty}^{\infty}\delta_\varepsilon^*(\phi)\,d\phi = 1$ を満たす関数である．
格子生成に適した具体的な定義として，ガウス型スムーズデルタ関数を用いる：

\begin{tcolorbox}[mybox, title={グリッド密度デルタ関数の定義（ガウス型）}]
\begin{equation}
  \delta_\varepsilon^*(\phi) =
  \frac{1}{\varepsilon_g\sqrt{\pi}}\exp\!\left(-\frac{\phi^2}{\varepsilon_g^2}\right)
  \label{eq:grid_delta}
\end{equation}
ここで $\varepsilon_g$ は界面近傍の格子細分化幅（有効半幅）を制御するパラメータである．

\textbf{正規化の確認：}
\[
  \int_{-\infty}^{\infty} \delta_\varepsilon^*(\phi)\,d\phi
  = \frac{1}{\varepsilon_g\sqrt{\pi}}\int_{-\infty}^{\infty}e^{-\phi^2/\varepsilon_g^2}\,d\phi
  = \frac{1}{\varepsilon_g\sqrt{\pi}} \cdot \varepsilon_g\sqrt{\pi} = 1 \quad \checkmark
\]

\textbf{パラメータの選択指針：}
\begin{itemize}
  \item $\varepsilon_g \approx 2\text{--}4 \, \varepsilon$（$\varepsilon$ は CLS 法の界面幅パラメータ）
  \item $\phi = 0$ では $\delta_\varepsilon^*(0) = 1/(\varepsilon_g\sqrt{\pi})$，界面での格子密度は $\omega(0) = \alpha$
  \item $|\phi| > 3\varepsilon_g$ の領域では $\delta_\varepsilon^*(\phi) < e^{-9}/(\varepsilon_g\sqrt{\pi}) \approx 0$：ほぼ一様格子に戻る
\end{itemize}
\end{tcolorbox}

\begin{tcolorbox}[warnbox, title={代替定義：コサイン型スムーズデルタ関数}]
ガウス型の代わりに，有限台を持つコサイン型を使うこともある（界面から完全に離れた点で厳密にゼロとなる利点）：
\[
  \delta_\varepsilon^*(\phi) =
  \begin{cases}
    \dfrac{1}{2\varepsilon_g}\left[1 + \cos\!\left(\dfrac{\pi\phi}{\varepsilon_g}\right)\right] & |\phi| \le \varepsilon_g \\[6pt]
    0 & |\phi| > \varepsilon_g
  \end{cases}
\]
この場合も $\int_{-\infty}^{\infty}\delta_\varepsilon^*\,d\phi = 1$ が満たされる（$[-\varepsilon_g,+\varepsilon_g]$ 上の積分が1）．
\end{tcolorbox}

$\alpha > 1$ は界面近傍の格子密度倍率である．

\subsection{座標変換とメトリクス係数}

一様計算空間 $\xi$ への座標変換を行う．物理空間微分と計算空間微分の関係は以下の通り厳密に導出される．

\begin{tcolorbox}[resultbox]
\textbf{座標変換公式（厳密形）：}
\begin{align}
 \pd{f}{x} &= J_x\,\pd{f}{\xi}
 \label{eq:transform_1st_correct} \\[6pt]
 \pdd{f}{x} &= J_x^2\,\pdd{f}{\xi} + J_x\,\pd{J_x}{\xi}\,\pd{f}{\xi}
 \label{eq:transform_2nd_correct}
\end{align}
\end{tcolorbox}
ここで $J_x = d\xi/dx$ はメトリクス係数であり，CCD 法を用いて高精度に評価される（次章 第\ref{sec:CCD}章で詳述）．

\subsection{格子点座標の生成アルゴリズム}
\label{sec:grid_gen}

メトリクス係数 $J_x = d\xi/dx$ を得るには，まず物理座標配列 $\{x_i\}$ を
グリッド密度関数 $\omega(\phi)$ から生成する必要がある．
以下に，台形公式を用いた累積和による具体的な手順を示す．

\begin{tcolorbox}[algbox, title={1次元格子点生成アルゴリズム（台形則積分）}]
\textbf{入力：} 格子数 $N$，領域長 $L$，初期水準集合値 $\{\phi^{(0)}_i\}$（一様格子上），密度倍率 $\alpha$

\textbf{ステップ1：} 一様計算格子 $\xi_i = iL/(N-1)$（$i=0,1,\ldots,N-1$）を定義する．

\textbf{ステップ2：} 各格子点で $\omega_i = \omega(\phi^{(0)}_i) = 1 + (\alpha-1)\delta_\varepsilon^*(\phi^{(0)}_i)$ を計算する．

\textbf{ステップ3：} 台形則で累積積分し，物理座標の「未正規化」値を求める：
\[
  \tilde{x}_0 = 0, \qquad
  \tilde{x}_{i+1} = \tilde{x}_i + \frac{\Delta\xi}{\omega_i},
  \quad \Delta\xi = \frac{L}{N-1}
\]

\textbf{ステップ4：} $\tilde{x}_{N-1}$ で割って領域長 $L$ に正規化する：
\[
  x_i = L \cdot \frac{\tilde{x}_i}{\tilde{x}_{N-1}}
\]

\textbf{ステップ5：} メトリクス係数を差分で求める：
\[
  J_x\big|_i = \frac{d\xi}{dx}\bigg|_i \approx \frac{\xi_{i+1}-\xi_{i-1}}{x_{i+1}-x_{i-1}}
  = \frac{2\Delta\xi}{x_{i+1}-x_{i-1}}
  \quad (\text{境界は片側差分})
\]

\textbf{出力：} 物理座標配列 $\{x_i\}$ とメトリクス係数 $\{J_x|_i\}$．
\end{tcolorbox}

\begin{tcolorbox}[warnbox, title={【重要】メトリクス係数 $J_x$ の精度と $\partial J_x/\partial\xi$ の評価}]
\label{box:grid_jx_accuracy}
上記の中心差分 $J_x|_i \approx 2\Delta\xi/(x_{i+1}-x_{i-1})$ は $\mathcal{O}(h^2)$ 精度しか持たない．
座標変換を用いた2階微分（式 \eqref{eq:transform_2nd_correct}）：
\[
  \pdd{f}{x} = J_x^2 \fpp{\xi} + J_x \pd{J_x}{\xi} \fp{\xi}
\]
の右辺第2項には $\partial J_x/\partial \xi$ が含まれる．
もし $J_x$ が $\mathcal{O}(h^2)$ 差分で与えられているとすると，
$\partial J_x/\partial\xi$ をさらに差分で評価すると精度は $\mathcal{O}(h)$ 以下に劣化し，
CCD の $\mathcal{O}(h^6)$ 空間精度が $\mathcal{O}(h)$ または $\mathcal{O}(h^2)$ に低下する．

\textbf{精度を保つための対処法（推奨）：}
\begin{enumerate}
  \item \textbf{$J_x$ を CCD で評価する：} 計算空間上で $J_x(\xi) = 1/(dx/d\xi)$ を CCD の $x'(\xi)$ として求め，$\mathcal{O}(h^6)$ 精度の $J_x$ と $\partial J_x/\partial\xi$ を得る．具体的には：
  \[
    J_x|_i = \frac{1}{(dx/d\xi)_i}, \qquad
    \pd{J_x}{\xi}\bigg|_i = -\frac{(d^2x/d\xi^2)_i}{(dx/d\xi)_i^2}
  \]
  ここで $(dx/d\xi)_i$ および $(d^2x/d\xi^2)_i$ は座標配列 $x(\xi)$ に CCD を適用して得る．
  \item \textbf{格子変化が緩やかな場合：} 界面近傍の格子集中が滑らかであれば（$\omega(\phi)$ が緩やかに変化），$J_x$ の $\mathcal{O}(h^2)$ 誤差が全体精度に与える影響は小さい場合がある．ただし，定量的確認（格子収束テスト）なしにこの仮定を採用すべきでない．
  \item \textbf{代替案：高次片側差分による $J_x$ 計算：} 4次以上の差分で $J_x$ を評価することで中間的な解決策となる（実装負荷は低い）．
\end{enumerate}
\textbf{本稿の実装選択：} 上記の方法1（CCD による $J_x$ 評価）を推奨する．
格子収束テスト（第\ref{sec:verification}章）で $\mathcal{O}(h^6)$ 収束が確認されれば，この評価が適切に機能していることを確認できる．
\end{tcolorbox}

\begin{tcolorbox}[defbox, title={アルゴリズムの物理的直感}]
$\omega(\phi) > 1$ となる界面近傍では $\Delta\tilde{x} = \Delta\xi/\omega < \Delta\xi$ となり，
物理格子幅が細くなる（$\omega$ 倍の格子密度）．
逆に界面から遠い領域（$\omega \approx 1$）では均一な格子幅が保たれる．
正規化（ステップ4）により，格子の総長が常に $L$ に一致することが保証される．
\end{tcolorbox}

\subsection{2次元格子生成への拡張}
\label{sec:grid_2d}

本節では，1次元の手順を2次元に拡張する方法を示す．
2次元では $x$ 方向と $y$ 方向の格子を\textbf{独立に}生成する（テンソル積格子）．

\subsubsection{テンソル積格子の生成手順}

\begin{tcolorbox}[algbox, title={2次元格子点生成アルゴリズム（テンソル積）}]
\textbf{前処理：} 初期 Level Set 場 $\phi^{(0)}_{i,j}$ を一様格子上で計算しておく．

\textbf{$x$ 方向格子：}
\begin{enumerate}
  \item $x$ 方向の密度関数 $\omega^x_i = 1 + (\alpha-1)\delta_\varepsilon^*(\bar{\phi}^x_i)$ を計算する．
        ここで $\bar{\phi}^x_i = \min_j |\phi^{(0)}_{i,j}|$（列 $i$ における界面までの最短距離）または
        $\bar{\phi}^x_i = (1/N_y)\sum_j \phi^{(0)}_{i,j}$（列平均）を用いる．
  \item 前節のステップ2--5で物理座標 $\{x_i\}$ とメトリクス係数 $\{J_x|_i\}$ を計算する．
\end{enumerate}

\textbf{$y$ 方向格子：}
\begin{enumerate}
  \item $y$ 方向の密度関数 $\omega^y_j = 1 + (\alpha-1)\delta_\varepsilon^*(\bar{\phi}^y_j)$ を計算する．
        ここで $\bar{\phi}^y_j = \min_i |\phi^{(0)}_{i,j}|$（行 $j$ における界面までの最短距離）を用いる．
  \item 同様の積分・正規化・メトリクス係数計算を実行し，$\{y_j\}$ と $\{J_y|_j\}$ を得る．
\end{enumerate}

\textbf{出力：} $N_x \times N_y$ の非一様格子 $\{(x_i, y_j)\}$ と
メトリクス係数 $\{J_x|_i\}$，$\{J_y|_j\}$（2次元に適用する際は $J_x$ は $i$ のみに，$J_y$ は $j$ のみに依存）．
\end{tcolorbox}

\begin{tcolorbox}[defbox, title={テンソル積格子の仮定と制限}]
テンソル積格子は「$x$ 方向格子と $y$ 方向格子の直積」であるため，以下の特性を持つ：
\begin{itemize}
  \item \textbf{直交格子：} 格子線は $x$ 方向と $y$ 方向に平行であり，常に直交する．
        変換行列のオフ対角項はゼロ（$\partial\xi/\partial y = \partial\eta/\partial x = 0$）．
  \item \textbf{界面傾き対応：} 界面が斜め（例：45°）の場合，界面近傍の格子細分化が
        $x$ 方向と $y$ 方向の両方で自動的に行われる（$\min_j|\phi_{i,j}|$ を使うため）．
  \item \textbf{制限：} 界面が1本の直線または単純閉曲線の場合には有効だが，
        複数の界面が互いに異なる方向を向く場合，一方の方向での細分化が過剰になる可能性がある．
\end{itemize}

\textbf{本稿での実装上の仮定：} 座標変換式（式 \eqref{eq:transform_1st_correct}, \eqref{eq:transform_2nd_correct}）は
直交テンソル積格子を前提としており，$\partial x/\partial\eta = \partial y/\partial\xi = 0$ が成立する．
非直交格子では追加のオフ対角メトリクス項が必要になるため，本稿の適用範囲外となる．
\end{tcolorbox}

\subsubsection{2次元格子の実装例（Python 擬似コード）}

以下の擬似コードは上記アルゴリズムの実装例を示す：

\begin{tcolorbox}[mybox, title={2次元格子生成の Python 擬似コード}]
\begin{verbatim}
import numpy as np

def build_1d_grid(phi_min_row, N, L, alpha, eps_g):
    """1次元格子生成（行または列の最小|phi|を入力）"""
    xi = np.linspace(0, L, N)
    dxi = L / (N - 1)
    # ガウス型スムーズデルタ関数
    omega = 1 + (alpha - 1) * np.exp(-phi_min_row**2 / eps_g**2) \
                             / (eps_g * np.sqrt(np.pi))
    # 台形則で累積積分
    x_tilde = np.zeros(N)
    for i in range(N - 1):
        x_tilde[i+1] = x_tilde[i] + dxi / omega[i]
    # 領域長 L に正規化
    x = L * x_tilde / x_tilde[-1]
    # メトリクス係数（中心差分）
    Jx = np.gradient(xi, x)  # dxi/dx
    return x, Jx

def build_2d_grid(phi0, Nx, Ny, Lx, Ly, alpha, eps_g):
    """2次元テンソル積格子生成"""
    # x 方向：各列の最小 |phi|
    phi_min_x = np.min(np.abs(phi0), axis=1)  # shape (Nx,)
    x, Jx = build_1d_grid(phi_min_x, Nx, Lx, alpha, eps_g)
    # y 方向：各行の最小 |phi|
    phi_min_y = np.min(np.abs(phi0), axis=0)  # shape (Ny,)
    y, Jy = build_1d_grid(phi_min_y, Ny, Ly, alpha, eps_g)
    # 2次元メッシュ
    X, Y = np.meshgrid(x, y, indexing='ij')
    return X, Y, Jx, Jy
\end{verbatim}
\end{tcolorbox}

\subsection{時間変化格子と ALE 効果}
\label{sec:grid_ale}

本稿の格子は初期水準集合値 $\phi^{(0)}$ から生成されるが，界面位置は時間発展する．
格子を毎タイムステップ更新する場合，\textbf{ALE（Arbitrary Lagrangian-Eulerian）定式化}が必要となる．

\begin{tcolorbox}[warnbox, title={本稿の実装方針：時間固定格子（簡単化のための仮定）}]
ALE 定式化では，格子速度 $\bm{v}_\text{mesh}$ を輸送方程式に追加する必要がある：
\[
  \frac{\partial\psi}{\partial t} + \nabla\cdot\bigl((\bm{u} - \bm{v}_\text{mesh})\psi\bigr) = 0
\]
ただし，この追加項の離散化は実装の複雑さを大きく増加させる．

\textbf{本稿の簡単化：} 格子を時間固定として扱う（$\bm{v}_\text{mesh} = 0$）か，
タイムステップごとに格子を再生成するが $\bm{v}_\text{mesh}$ の補正は省略する．
\end{tcolorbox}

\begin{tcolorbox}[defbox, title={ALE 省略による精度劣化の定量的見積もり}]
ALE 補正項 $\nabla\cdot(\bm{v}_\text{mesh}\psi)$ を省略したときの誤差を定量的に評価する．

\textbf{誤差の発生機構：}
1ステップあたりの界面移動量を $\delta = |\bm{u}|\Delta t$ とすると，
格子が再生成される場合，界面近傍の格子点は $\delta/h$ 程度移動する（$h$：格子間隔）．
ALE 補正を省略することで，$\bm{v}_\text{mesh}$ による追加対流量 $\bm{v}_\text{mesh}\cdot\bnabla\psi \sim |\bm{v}_\text{mesh}| \cdot |\bnabla\psi|$ が無視される．

\textbf{界面移動速度と格子速度の関係：}
界面集中格子では $|\bm{v}_\text{mesh}| \approx \dot{\bm{x}}_{\text{interface}} \approx |\bm{u}|_{\Gamma}$（界面速度）である．
これを省略した場合，1ステップ当たりの誤差は：
\[
  \varepsilon_{\text{ALE}} \sim |\bm{v}_\text{mesh}| \cdot |\bnabla\psi| \cdot \Delta t \sim U_\Gamma \cdot \frac{1}{\varepsilon} \cdot \Delta t
\]
ここで $|\bnabla\psi| \sim 1/(4\varepsilon)$（界面近傍のプロファイル），$U_\Gamma$ は界面速度である．

\textbf{近似が有効な条件：}
\begin{center}
\renewcommand{\arraystretch}{1.3}
\begin{tabular}{>{\raggedright\arraybackslash}p{48mm}c>{\raggedright\arraybackslash}p{42mm}}
\hline
条件 & 移動量 $\delta/h$ & ALE 省略の妥当性 \\
\hline
緩やかな界面変形 & $< 0.1$ & 良好（誤差 $< 10\%$ of $\Delta t^2$ 項） \\
標準的な問題 & $\approx 0.5$ & 条件付きで可（格子収束確認要） \\
大密度比 ($\rho_l/\rho_g{=}1000$) の高速界面 & $> 1$ & 精度劣化の可能性大（ALE 実装推奨） \\
\hline
\end{tabular}
\end{center}

\textbf{大密度比問題での注意：}
$\rho_l/\rho_g = 1000$ の場合，液相変形が速く $\delta/h > 1$ となることがある．
この状況では ALE 補正なしの格子更新は実質的に $O(\Delta t)$ 精度（1次）となり，
時間積分スキーム（TVD-RK3, Crank-Nicolson の $O(\Delta t^3), O(\Delta t^2)$）の精度優位性が失われる．

\textbf{推奨する検証手順：}
格子収束テスト（第\ref{sec:verification}章）に加えて，時間刻み $\Delta t$ を半分にしたとき誤差が $1/4$ 倍になる（2次収束）か，
$1/8$ 倍（3次収束）になるかを確認し，ALE 省略の影響を定量的に評価する．
\end{tcolorbox}
